\documentclass[UTF8,11pt,a4paper]{ctexart} 
\usepackage{graphicx} 
\usepackage{float}
\usepackage{geometry}
\usepackage{xcolor}
\usepackage{fancyhdr}
\includeonly{loser_empire.tex}
\geometry{left=2cm,right=2cm,top=3cm,bottom=2cm}
\pagestyle{fancy}
\lhead{\author}
\chead{\date}
\rhead{1 你好卢瑟}
\lfoot{}
\cfoot{\thepage}
\rfoot{}
\renewcommand{\headrulewidth}{0.2pt}
\renewcommand{\headwidth}{\textwidth}
\renewcommand{\footrulewidth}{0pt}
\title{卢瑟帝国宪法 }
\author{mq白}
\date{\today}
\begin{document}
\maketitle
\section{你好 卢瑟}
事实上在阅读宪法之前，请先明确自己的身份，是“\textbf{卢瑟}”，否则则不需要遵守。
不管什么职位，都是 \textbf{卢瑟}。
\subsection{你好 卢瑟王}
\paragraph{卢瑟王}是卢瑟帝国的王者，可以统御普通卢瑟。
\subsection{底层卢瑟你好}
\paragraph{国王}在卢瑟帝国有统治权。\textbf{mq白}是国王，管理一切， 是卢瑟帝国中心。
\subsection{你好 卢瑟管理者}
\paragraph{卢瑟管理者}是专门管理卢瑟的，地位和卢瑟王一样。

\subsection{你好 赞助商}
\paragraph{赞助商}是卢瑟帝国唯一的资金来源，非常重要，用于供养国王和帝国日常发展。\\
{\color{white}-}
\begin{figure}[H]
    \caption{简单等级示意图} 
    \centering 
    \includegraphics[width = .7\textwidth]{resource/luse.png}
\end{figure}

\begin{center}
    \begin{tabular}{|l|c|r|}
        \hline
             & \textbf{卢瑟帝国国王} &       \\
        \hline
        卢瑟王  & 赞助商             & 卢瑟管理者 \\
        \hline
        mq家族 & 史官              & 底层卢瑟  \\
        \hline
             & 新子民             &       \\
        \hline
    \end{tabular}
\end{center} 

\end{document}